
\textbf{Remarque 4.9.} Dans le cas @@M0@@ noethérien, @@M1@@ régulier, on peut, en procédant comme en 4.7, décrire la structure d'anneau graduée de @@M2@@ en termes des anneaux gradués @@M3@@, @@M4@@, des homomorphismes additifs @@M5@@ et @@M6@@ déduits de @@M7@@ et @@M8@@, et de l'élément @@M9@@ de @@M10@@. De même, on peut décrire à partir des mêmes données les homomorphismes @@M11@@ définis à partir de @@M12@@ par passage aux gradués associés.

\subsection*{4.10.}

Soit @@M13@@ un schéma \emph{régulier} ayant un faisceau ample. On sait (IV\qquad) que les groupes de Grothendieck formés avec les @@M14@@-Modules localement libres de type fini, et avec les @@M15@@-Modules cohérents, sont isomorphes. On utilise ce fait pour définir sur @@M16@@ une filtration « topologique », en prenant pour @@M17@@ le sous-groupe engendré par les classes des faisceaux cohérents dont le support est de codimension @@M18@@ ; on montre facilement que @@M19@@ est engendré par les classes des faisceaux @@M20@@, où @@M21@@ parcourt l'ensemble des sous-schémas fermés intègres de @@M22@@, de codimension @@M23@@(X\qquad).

\textbf{Proposition 4.11.} \uline{Soit} @@M24@@ \uline{un schéma régulier ayant un faisceau ample}. \uline{Alors} \uline{pour tout} @@M25@@ :
@@M26@@
@@M27@@ L'inclusion inverse, valable sans tensorisation par @@M28@@, sera prouvée dans X\qquad).

D'après 4.10, il suffit de montrer que si @@M29@@ est un sous-schéma fermé intègre de @@M30@@, de codimension @@M31@@, alors :
@@M32@@

Pour cela, on procède par récurrence noethérienne sur les sous-schémas fermés intègres de @@M33@@. Supposons donc la propriété prouvée pour les sous-schémas fermés intègres de @@M34@@, et montrons-la pour @@M35@@. Soit @@M36@@ l'Idéal de @@M37@@
